\begin{abstract}
Crop-ranking systems order crops by agronomic or predicted performance, whereas
acreage decisions depend on uncertain margins, joint downside risk and
operational feasibility. We formulate a stochastic crop-planning programme
with price, yield and cost uncertainty; land, budget, rotation, contract,
labour and equipment constraints; and a loss-CVaR ceiling. In a fixed-land
two-crop model, when the higher-ranked crop also has the higher expected margin,
reversal occurs if and only if the risk-feasible share interval ends below the
equal-share allocation. In a
public-data-calibrated Kansas experiment, winter wheat ranks first by a
pre-decision historical relative-yield performance index but receives the
smallest optimized share, producing a complete rank reversal. The Kansas
result is margin-induced and moderated by downside risk. Separate controlled
experiments show that tightening a loss-CVaR ceiling can
reverse the soybean--corn allocation order while preserving both score and
expected-margin order, whereas a soybean rotation cap produces an
operationally induced reversal. A selected Gaussian mean--variance policy
reduces Gaussian variance by \DiversificationVarianceReduction\% relative to
the expected-profit benchmark but violates the Student-\(t\) evaluation-law
loss-CVaR ceiling. Under a shared ex-ante loss-CVaR model, the
information--flexibility interaction can be positive, zero or negative,
depending on whether flexibility enables reallocation or buffers the predicted
state. A \StateYears{}-state-year public panel documents score--acreage
disagreement descriptively, while its lagged association includes zero and
does not identify private objectives or constraints. Crop rankings therefore
support acreage recommendations only when combined with explicit economic,
risk and operational assumptions.
\end{abstract}
